% ============================================================
% ODE IVP Project --- LaTeX Report
% Compile: pdflatex report.tex  (twice for the TOC)
% Structure follows the project brief:
%   Background -> Model -> Methods -> Stability -> Experiments -> Conclusions
%
% This is the MERGED template: it combines the project-specific guidance for
% Project 1 (ODE IVPs) with the course-level report skeleton (algorithm
% environments, theorem environments, convergence table, AI Transparency Log
% and ICS appendices).
%
% Project 1, Direction 4: Financial SDE --- GBM benchmark and non-affine
% stochastic volatility.  The deterministic stability/adaptivity requirements
% of the ODE directions are replaced by exact--EM--Milstein comparison,
% strong/weak convergence with coupled paths, Monte Carlo uncertainty, and
% nested-grid validation for the non-affine model.
% ============================================================

\documentclass[11pt,a4paper]{article}

\usepackage[utf8]{inputenc}
\usepackage[T1]{fontenc}
\usepackage{amsmath, amssymb, amsthm}
\usepackage{geometry}
\geometry{a4paper, margin=2.5cm}
\usepackage{graphicx}
\usepackage{booktabs}
\usepackage{listings}
\usepackage{xcolor}
\usepackage{hyperref}
\usepackage{enumitem}
\usepackage{float}
\usepackage{caption}
\usepackage{subcaption}
\usepackage{bm}
\usepackage{algorithm}
\usepackage{algpseudocode}
\usepackage{mdframed}

% --- colors ---
\definecolor{lecture}{RGB}{70,130,180}
\definecolor{seminar}{RGB}{205,92,92}
\definecolor{formbg}{RGB}{250,250,245}
\definecolor{formframe}{RGB}{180,170,140}

% --- code listing style (Python) ---
\lstset{
    language=Python,
    basicstyle=\ttfamily\scriptsize,
    keywordstyle=\color{blue}\bfseries,
    commentstyle=\color{green!40!black},
    stringstyle=\color{red!70!black},
    numbers=left,
    numberstyle=\tiny\color{gray},
    frame=single,
    breaklines=true,
    showstringspaces=false,
    tabsize=4,
    captionpos=b,
}

% --- theorem-like environments (course-level set) ---
\newtheorem{theorem}{Theorem}[section]
\newtheorem{lemma}[theorem]{Lemma}
\newtheorem{proposition}[theorem]{Proposition}
\newtheorem{corollary}[theorem]{Corollary}
\theoremstyle{definition}
\newtheorem{definition}[theorem]{Definition}
\newtheorem{assumption}[theorem]{Assumption}
\newtheorem{example}[theorem]{Example}
\theoremstyle{remark}
\newtheorem{remark}[theorem]{Remark}

% --- algorithm style ---
\algrenewcommand{\algorithmicrequire}{\textbf{Input:}}
\algrenewcommand{\algorithmicensure}{\textbf{Output:}}
\algrenewcommand{\algorithmiccomment}[1]{\hfill\textcolor{gray}{\# #1}}

% --- form box (for appendices) ---
\newmdenv[
    backgroundcolor=formbg,
    linecolor=formframe,
    linewidth=0.8pt,
    roundcorner=3pt,
    innermargin=5pt,
    outermargin=5pt,
    innertopmargin=8pt,
    innerbottommargin=8pt
]{formbox}

\title{\textbf{Strong and Weak Convergence of Euler--Maruyama and Milstein
Schemes for Geometric Brownian Motion and a Non-Affine Stochastic Volatility
Model}\\[2mm]
\large Project 1 (ODE IVP) --- Direction 4: Financial SDE}
\author{
    Team \#\texttt{<N>} \\
    \small \texttt{<name1>, <id1> (Project Manager)} \quad
    \texttt{<name2>, <id2> (Mathematical Theory)} \\
    \small \texttt{<name3>, <id3> (Algorithm Implementation)} \quad
    \texttt{<name4>, <id4> (Visualization \& Report)} \\
    \small \texttt{<name5>, <id5> (Testing \& Validation)} \\
    \small Department of Applied Mathematics
}
\date{\today}

\begin{document}
\maketitle

% ============================================================
\begin{abstract}
\noindent
[150--250 words answering: (1) what problem and why interesting --- state that
two models are compared, geometric Brownian motion as an exact benchmark and a
non-affine stochastic volatility model with no closed-form transition law;
(2) what methods you implemented --- Euler--Maruyama and Milstein on the price
process, Euler--Maruyama on the two-dimensional non-affine state, and a
Brownian engine that produces one path shared by all step sizes; (3) key
findings --- the measured strong orders, the weak orders, whether the two weak
curves coincide, and whether the increments preserve their moments under
coarsening; (4) the main limitation --- that the non-affine reference is
numerical, so its rate is empirical, and that no scalar Milstein result applies
there.]
\end{abstract}

\noindent\textbf{Keywords:} stochastic differential equations,
Euler--Maruyama, Milstein method, strong convergence, weak convergence,
mean-square stability, It\^o--Taylor expansion, non-affine diffusion,
Monte Carlo uncertainty, geometric Brownian motion

\tableofcontents
\newpage

% ============================================================
\section{Introduction and Motivation}
% ============================================================

\subsection{Problem Background}

\subsubsection{SDEs in Finance}

A deterministic ordinary differential equation produces a single future for a
single initial condition. Asset prices do not behave in this way: two investors
who hold the same stock on the same day will not see the same price path
tomorrow. Randomness therefore has to enter the model itself, and the standard
way to do this is to add a Brownian term to the differential equation. The
result is a stochastic differential equation (SDE) of the general form
\begin{equation}
    dX_t = a(X_t, t)\,dt + b(X_t, t)\,dW_t,
    \label{eq:general-sde}
\end{equation}
where $a$ is the drift, $b$ is the diffusion coefficient, and $W_t$ is a
standard Brownian motion. The two coefficients play different roles. The drift
describes the average tendency of the process, while the diffusion describes
how strongly random shocks are amplified. Their relative size determines how
much of the observed variation is systematic and how much is noise.

The classical starting point is the geometric Brownian motion (GBM) model of
Black, Merton and Scholes~\cite{blackscholes}, in which both the drift and the
diffusion are proportional to the current price:
\begin{equation}
    dS_t = \mu S_t\,dt + \sigma S_t\,dW_t .
    \label{eq:gbm}
\end{equation}
Because of its exponential structure, \eqref{eq:gbm} can be solved in closed
form, and its transition law is known exactly. This makes GBM an ideal
\emph{verification} problem: any code written for it can be compared against a
reference that contains no discretisation error at all.

The same property, however, makes GBM unsuitable as a \emph{modelling} problem.
Constant volatility cannot reproduce two features that are easy to observe in
real option markets: the implied volatility smile and the tendency of
volatility to rise when prices fall. Modelling these effects requires the
volatility itself to be a random process, which leads to the second model
studied here. The report therefore uses two models with deliberately different
status: GBM supplies an exact benchmark, while the stochastic volatility model
is a problem in which numerical integration is genuinely part of the solution
rather than a check of a known answer.

\subsubsection{It\^o's Lemma}
\label{sec:ito}

Equation \eqref{eq:gbm} is written in the price $S_t$, but it is more
convenient to work with the log-price $X_t = \log S_t$. The chain rule of
ordinary calculus cannot be used for this change of variable, because a
Brownian path has infinite variation. The correct tool is It\^o's lemma: for
$f \in C^{1,2}$ and $X_t$ satisfying \eqref{eq:general-sde},
\begin{equation}
    df(X_t, t) =
    \left( \frac{\partial f}{\partial t}
         + a \frac{\partial f}{\partial x}
         + \tfrac12 b^2 \frac{\partial^2 f}{\partial x^2} \right) dt
    + b \frac{\partial f}{\partial x}\, dW_t .
    \label{eq:ito}
\end{equation}
The extra second-order term in \eqref{eq:ito} is the essential difference
between stochastic and deterministic calculus.

Applying \eqref{eq:ito} with $f(s) = \log s$, $a = \mu S_t$ and
$b = \sigma S_t$ gives
\begin{equation}
    dX_t = \left( \mu - \tfrac12 \sigma^2 \right) dt + \sigma\, dW_t ,
    \label{eq:logprice}
\end{equation}
so the log-price satisfies an SDE with \emph{constant} coefficients. The
$-\tfrac12\sigma^2$ term is not a modelling choice; it is forced by It\^o's
lemma, and dropping it changes the expected value of the price by a factor of
$e^{\sigma^2 T/2}$. Integrating \eqref{eq:logprice} gives the exact solution
\begin{equation}
    S(T) = S_0 \exp\left[ \left( \mu - \tfrac12 \sigma^2 \right) T
                        + \sigma W(T) \right] .
    \label{eq:gbm-exact}
\end{equation}

Two consequences of \eqref{eq:logprice} shape the whole numerical study, and
they are made precise in Section~\ref{sec:gbm-model}. Since the drift and the
diffusion are constant, the transition law of $X$ is exactly Gaussian, so the
updating rule applied to $X$ reproduces the exact solution at every grid point.
At the same time, the constant diffusion means the noise is additive: the
Milstein correction, which is proportional to the derivative of the diffusion
coefficient, vanishes identically on $X$, and no strong order can be measured
there. To obtain distinct orders, the price $S$ itself must be discretised,
where the noise $\sigma S$ does depend on the state.

We finally note the convention issue. The It\^o form is used throughout, which
is the standard convention in finance and is the reason the $-\tfrac12\sigma^2$
correction appears above. The same convention produces the
$-\tfrac12 g(Y_t)^2$ term in the stochastic volatility model introduced next.

\subsubsection{Non-Affine Stochastic Volatility}
\label{sec:non-affine}

The observation that motivates a stochastic volatility model is simple:
implied volatility varies with strike and maturity, and it tends to increase
when the underlying price falls. A constant $\sigma$ cannot generate either
pattern. The most widely used remedy is the Heston model~\cite{heston}, in
which the variance $V_t$ follows a mean-reverting square-root diffusion.
Heston belongs to the class of \emph{affine} diffusions: both its drift and its
covariance are affine functions of the state, and consequently the
characteristic function satisfies a Riccati ordinary differential equation that
can be solved in closed form. Prices are then recovered by Fourier inversion.
Affine structure is therefore not a technical detail; it is exactly what makes
a model analytically tractable~\cite{dfs}.

The model used in this report keeps the economic motivation but gives up that
tractability. Writing $X_t = \log S_t$, its two equations are
\[
    dX_t = \left( \mu - \tfrac12 g(Y_t)^2 \right) dt + g(Y_t)\, dW_t^{(1)},
    \qquad
    dY_t = \kappa (\theta - Y_t)\, dt + \xi \sqrt{1 + Y_t^{2}}\; dW_t^{(2)},
\]
with a bounded logistic volatility function $g$ and correlated Brownian
motions. The factor $Y_t$ is mean-reverting, so volatility is persistent rather
than independent from period to period, and a negative correlation produces the
leverage effect described above. The logistic function keeps the instantaneous
volatility strictly between two positive bounds, which removes the
vanishing- and negative-variance pathologies that arise in square-root models.
Simulating $X$ rather than $S$ keeps the price positive by construction. The
system is stated formally, with all parameter values, in
Section~\ref{sec:nasv-model}.

What makes this model \emph{non-affine} is visible in the two coefficients. The
diffusion of $Y_t$ is $\xi\sqrt{1 + Y_t^{2}}$, which is not an affine function
of $Y_t$, and the drift of $X_t$ depends on the square of a logistic function.
Both facts break the affine structure described above. The model therefore has
no closed-form transition density, no closed-form characteristic function, and
no practical exact sampler. This distinction should be stated precisely: the
equations do have an explicit integral representation, but it is not a usable
simulation formula. Time discretisation is consequently part of the model
solution, and the numerical reference against which coarse schemes are measured
must itself be validated by refinement, rather than taken from a formula.

\subsection{Mathematical Models}
\label{sec:models}

This section fixes the notation and states the two models precisely.
Section~\ref{sec:bm} defines the driving noise, Section~\ref{sec:gbm-model}
states the GBM benchmark, and Section~\ref{sec:nasv-model} states the
non-affine stochastic volatility model. All parameter values given here are the
values used by the code described in Section~3, and both models are posed on
the interval $0 \le t \le T = 1$.

\subsubsection{Brownian Motion}
\label{sec:bm}

\begin{definition}[Standard Brownian motion]
\label{def:bm}
A real-valued process $\{W_t\}_{t \ge 0}$ on a filtered probability space
$(\Omega, \mathcal{F}, (\mathcal{F}_t)_{t \ge 0}, \mathbb{P})$ is a
\emph{standard Brownian motion} if
\begin{enumerate}[label=(\roman*), leftmargin=2em]
    \item $W_0 = 0$ almost surely;
    \item for every $0 \le s < t$ the increment $W_t - W_s$ is independent of
          $\mathcal{F}_s$ and $W_t - W_s \sim \mathcal{N}(0, t - s)$;
    \item the map $t \mapsto W_t$ is continuous almost surely.
\end{enumerate}
\end{definition}

Three properties of this process are used repeatedly below. First,
$\mathbb{E}[W_t] = 0$ and $\operatorname{Var}[W_t] = t$, and $\{W_t\}$ is a
martingale. Second, the quadratic variation of a Brownian path is
$[W]_t = t$, which is why the ordinary chain rule fails and the second-order
term in It\^o's lemma appears. Third, Brownian paths are continuous but nowhere
differentiable and have infinite total variation on every interval, so
\eqref{eq:logprice} cannot be read as a classical differential equation.

To couple two sources of randomness, we use a two-dimensional Brownian motion
$\bm{W}_t = (W_t^{(1)}, W_t^{(2)})$ with
\begin{equation}
    d\langle W^{(1)}, W^{(2)} \rangle_t = \rho\, dt,
    \qquad |\rho| < 1 .
    \label{eq:correlation}
\end{equation}
Such a pair can always be built from two independent standard Brownian motions
$B^{(1)}$ and $B^{(2)}$ by setting
$W^{(1)} = B^{(1)}$ and $W^{(2)} = \rho B^{(1)} + \sqrt{1 - \rho^2}\, B^{(2)}$,
which is the Cholesky factorisation of the correlation matrix. On a uniform
grid $t_n = n h$ with $h = T / N$, this gives the increment convention used
throughout the code:
\begin{equation}
    \Delta W_n^{(1)} = \sqrt{h}\, Z_{1,n},
    \qquad
    \Delta W_n^{(2)} = \sqrt{h}\left( \rho Z_{1,n}
                       + \sqrt{1 - \rho^2}\, Z_{2,n} \right),
    \label{eq:increments}
\end{equation}
with $Z_{1,n}, Z_{2,n} \sim \mathcal{N}(0, 1)$ independent. Each increment then
has variance $h$, and
$\operatorname{Cov}\!\left( \Delta W_n^{(1)}, \Delta W_n^{(2)} \right) = \rho h$.
The symbol $h$ is used for the step size in this report; the corresponding
variable in the code is named \texttt{dt}.

\subsubsection{Geometric Brownian Motion}
\label{sec:gbm-model}

\paragraph{Notation and Parameters}
The single state variable is the price $S_t \in (0, \infty)$ itself, so the
state space has dimension one. The problem pack does not prescribe GBM
parameters, since GBM serves here only as a verification benchmark; the values
below are therefore chosen explicitly and used in every GBM experiment.

\begin{table}[H]
\centering
\begin{tabular}{llll}
\toprule
\textbf{Symbol} & \textbf{Meaning} & \textbf{Value} & \textbf{Unit} \\
\midrule
$S_0$    & initial price        & $100$   & --- \\
$\mu$    & drift rate           & $0.05$  & yr$^{-1}$ \\
$\sigma$ & volatility           & $0.20$  & yr$^{-1/2}$ \\
$T$      & final time           & $1$     & yr \\
$h$      & step size, $h = T/N$ & varies  & yr \\
\bottomrule
\end{tabular}
\caption{GBM benchmark parameters and step size.}
\label{tab:gbm-params}
\end{table}

\paragraph{SDE System}
\begin{equation}
    dS_t = \mu S_t\, dt + \sigma S_t\, dW_t,
    \qquad S_0 = 100,
    \qquad 0 \le t \le T = 1 .
    \label{eq:gbm-model}
\end{equation}
Both coefficients are proportional to $S_t$, so they are globally Lipschitz and
grow at most linearly. The model is autonomous: neither coefficient depends
explicitly on $t$.

\paragraph{Exact Solution, Transition and Moments}
The continuous-time solution was derived in Section~\ref{sec:ito} and is given
by \eqref{eq:gbm-exact}. Because the drift and the diffusion of the log-price
\eqref{eq:logprice} are constant, the transition law over one step is exactly
Gaussian:
\begin{equation}
    X_{n+1} \mid X_n \;\sim\;
    \mathcal{N}\!\left( X_n + \left( \mu - \tfrac12 \sigma^2 \right) h,\;
                    \sigma^2 h \right),
    \qquad X_n = \log S_n .
    \label{eq:gbm-transition}
\end{equation}
Sampling this transition gives the exact simulation scheme
\begin{equation}
    X_{n+1} = X_n + \left( \mu - \tfrac12 \sigma^2 \right) h
              + \sigma \sqrt{h}\, Z_n,
    \qquad
    S_{n+1} = \exp\left( X_{n+1} \right),
    \label{eq:gbm-sampler}
\end{equation}
with $Z_n \sim \mathcal{N}(0,1)$ independent. Iterating \eqref{eq:gbm-sampler}
is equivalent to applying Euler--Maruyama to the log-price equation and
exponentiating: the two are identical at every grid point, so
\eqref{eq:gbm-sampler} carries \emph{no discretisation error}. This is what
makes it the reference in the strong convergence study of
Section~\ref{sec:experiments}, and it is also why no convergence order can be
measured on the log-price scale.

The exact first and second moments follow directly from
\eqref{eq:gbm-exact}:
\begin{equation}
    \mathbb{E}[S(T)] = S_0 e^{\mu T} = 105.1271,
    \qquad
    \operatorname{Var}[S(T)] = S_0^2 e^{2\mu T}
        \left( e^{\sigma^2 T} - 1 \right) = 451.03 .
    \label{eq:gbm-moments}
\end{equation}
The corresponding standard deviation is $21.24$, which sets the Monte Carlo
sample size needed to resolve a given bias in Section~\ref{sec:experiments}.
The distribution of $S(T)$ is lognormal with parameters
$\left( \log S_0 + \left( \mu - \tfrac12\sigma^2 \right) T,\; \sigma^2 T
\right)$, and this is the law against which the simulated distribution is
compared.

\subsubsection{Non-Affine Stochastic Volatility}
\label{sec:nasv-model}

\paragraph{Notation and Parameters}
The state is the pair $z_t = (X_t, Y_t) \in \mathbb{R}^2$, where
$X_t = \log S_t$ is the log-price and $Y_t$ is a dimensionless volatility
factor. The price is recovered as $S_t = e^{X_t}$, so $S_t > 0$ holds by
construction. The parameters are those of the reproducible benchmark in the
problem pack.

\begin{table}[H]
\centering
\begin{tabular}{llll}
\toprule
\textbf{Symbol} & \textbf{Meaning} & \textbf{Value} & \textbf{Unit} \\
\midrule
$S_0$                         & initial price                  & $100$  & --- \\
$Y_0$                         & initial volatility factor      & $0$    & --- \\
$\mu$                         & drift rate of the log-price    & $0.05$ & yr$^{-1}$ \\
$\kappa$                      & mean-reversion speed of $Y$    & $2$    & yr$^{-1}$ \\
$\theta$                      & long-run level of $Y$          & $-0.2$ & --- \\
$\xi$                         & volatility-of-volatility       & $0.6$  & yr$^{-1/2}$ \\
$\rho$                        & correlation of $W^{(1)}, W^{(2)}$ & $-0.7$ & --- \\
$\sigma_{\min}$               & lower bound of $g$             & $0.10$ & yr$^{-1/2}$ \\
$\sigma_{\max}$               & upper bound of $g$             & $0.50$ & yr$^{-1/2}$ \\
$T$                           & final time                     & $1$    & yr \\
\bottomrule
\end{tabular}
\caption{Non-affine stochastic volatility benchmark parameters.}
\label{tab:nasv-params}
\end{table}

Two derived quantities are useful for orientation. Since $g$ is logistic,
$g(0) = (\sigma_{\min} + \sigma_{\max})/2 = 0.30$ is the instantaneous
volatility at $t = 0$, while $g(\theta) = 0.2801$ is the level that volatility
reverts to. Because $Y_0 = 0 > \theta$, the volatility factor starts above its
long-run level and drifts downwards over the year.

\paragraph{SDE System}
\begin{align}
    dX_t &= \left( \mu - \tfrac12 g(Y_t)^2 \right) dt
            + g(Y_t)\, dW_t^{(1)},
    \label{eq:nasv-model-x} \\[2pt]
    dY_t &= \kappa (\theta - Y_t)\, dt
            + \xi \sqrt{1 + Y_t^{2}}\; dW_t^{(2)},
    \label{eq:nasv-model-y}
\end{align}
with the bounded logistic volatility function
\begin{equation}
    g(y) = \sigma_{\min} + \frac{\sigma_{\max} - \sigma_{\min}}{1 + e^{-y}},
    \qquad 0 < \sigma_{\min} < \sigma_{\max},
    \label{eq:nasv-g}
\end{equation}
and the correlation \eqref{eq:correlation}. The initial condition is
$(X_0, Y_0) = (\log 100,\, 0)$ and the interval is again $0 \le t \le T = 1$.
The system is autonomous. The drift of $X$ depends on $Y$ only, so $X_t$ does
not feed back into $Y_t$; the pair is one-way coupled.

\begin{assumption}[Regularity of the coefficients]
\label{ass:regularity}
For both systems \eqref{eq:gbm-model} and \eqref{eq:nasv-model-x}--%
\eqref{eq:nasv-model-y}, the drift and diffusion coefficients are globally
Lipschitz and satisfy a linear growth bound.
\end{assumption}

Assumption~\ref{ass:regularity} is easy to verify for these two models, and the
constants are small. For the logistic function,
$g'(y) \le (\sigma_{\max} - \sigma_{\min})/4 = 0.10$, so $g$ is Lipschitz with
constant $0.10$ and bounded by $0.50$. The drift of $X$ depends on $Y$ through
$\tfrac12 g^2$, whose derivative is at most $g g' \le 0.05$. The diffusion of
$Y$ satisfies $\left| \xi \sqrt{1 + y^2} \right| \le \xi (1 + |y|)$, so it grows
linearly with constant $\xi = 0.6$ and is Lipschitz with the same constant.
Because these bounds hold uniformly, a unique strong solution of each system
exists for all $t \ge 0$. Assumption~\ref{ass:regularity} is exactly the
hypothesis invoked by the convergence statements in
Section~\ref{sec:convergence}, so it is stated here once and referenced there.

\begin{remark}[An exact identity that survives]
\label{rem:nasv-mean}
The boundedness of $g$ gives one exact result even though the model is not
tractable. Writing the price equation in the form
$dS_t = \mu S_t\, dt + g(Y_t) S_t\, dW_t^{(1)}$, the discounted price
$S_t e^{-\mu t}$ is a local martingale whose diffusion coefficient is bounded.
Novikov's condition therefore holds, since
$\tfrac12 \int_0^T g(Y_u)^2\, du \le \tfrac12 \sigma_{\max}^2 T = 0.125$ is a
deterministic bound, and consequently
\begin{equation}
    \mathbb{E}[S_T] = S_0 e^{\mu T} = 105.1271 .
    \label{eq:nasv-mean}
\end{equation}
This identity is exact and independent of $\kappa, \theta, \xi, \rho,
\sigma_{\min}$ and $\sigma_{\max}$. It is used in
Section~\ref{sec:impl-verify} as a hard test of the implementation, because an
error in the $-\tfrac12 g(Y_t)^2$ term of \eqref{eq:nasv-model-x} would shift
the sample mean away from $105.1271$ by a detectable amount.
\end{remark}

What the model does \emph{not} provide is an analytically usable transition
law. As explained in Section~\ref{sec:non-affine}, the quadratic dependence of
the variance on $Y_t$ destroys the affine structure, so there is no closed-form
transition density, no closed-form characteristic function and no practical
exact sampler. Discretising \eqref{eq:nasv-model-x}--\eqref{eq:nasv-model-y} is
therefore part of the solution of the model, and the reference used in
Section~\ref{sec:experiments} is a finely resolved numerical solution rather
than a formula.

\subsection{Overview of the Report}

Section~\ref{sec:discretisation} derives the two discrete schemes, explains
where their convergence orders come from, and discusses which of the usual ODE
tools do and do not transfer to SDEs. Section~\ref{sec:implementation}
describes the implementation, the shared-path Brownian engine and the checks
applied to the code before any result is quoted.
Section~\ref{sec:experiments} presents the numerical experiments, and
Section~\ref{sec:conclusions} concludes.

\newpage
% ============================================================
\section{Discretisation and Analysis}
\label{sec:discretisation}
% ============================================================

This section derives the two schemes used for GBM and the scheme used for the
non-affine model, explains where their convergence orders come from, and states
clearly which of the usual ODE tools do and do not transfer to SDEs. All
statements below refer to the models of Section~\ref{sec:gbm-model} and
Section~\ref{sec:nasv-model}, and to Assumption~\ref{ass:regularity}.

\subsection{Discrete Schemes}
\label{sec:schemes}

\subsubsection{Euler--Maruyama}

Consider the general scalar It\^o SDE $dX_t = a(X_t)\,dt + b(X_t)\,dW_t$.
Integrating over one step $[t_n, t_{n+1}]$ with $h = t_{n+1} - t_n$ gives
\begin{equation}
    X_{n+1} = X_n + \int_{t_n}^{t_{n+1}} a(X_s)\, ds
                   + \int_{t_n}^{t_{n+1}} b(X_s)\, dW_s .
    \label{eq:integral-form}
\end{equation}
The Euler--Maruyama (EM) scheme is obtained by freezing both integrands at the
left endpoint:
\begin{equation}
    X_{n+1} = X_n + a(X_n)\, h + b(X_n)\, \Delta W_n,
    \qquad \Delta W_n = W_{t_{n+1}} - W_{t_n} .
    \label{eq:em-general}
\end{equation}
The deterministic integral is approximated by a rectangle rule and the
It\^o integral by its integrand at the left endpoint; both approximations are
consistent to first order.

Applied to the two models of this report, \eqref{eq:em-general} gives
\begin{align}
    \text{(GBM, price)} \quad
    S_{n+1} &= S_n + \mu S_n h + \sigma S_n \Delta W_n,
    \label{eq:em-gbm} \\[2pt]
    \text{(GBM, log-price)} \quad
    X_{n+1} &= X_n + \left( \mu - \tfrac12 \sigma^2 \right) h
               + \sigma \Delta W_n,
    \label{eq:em-gbm-log} \\[2pt]
    \text{(non-affine SV)} \quad
    X_{n+1} &= X_n + \left( \mu - \tfrac12 g(Y_n)^2 \right) h
               + g(Y_n) \Delta W_n^{(1)}, \notag \\[2pt]
    Y_{n+1} &= Y_n + \kappa (\theta - Y_n) h
               + \xi \sqrt{1 + Y_n^{2}}\, \Delta W_n^{(2)} .
    \label{eq:em-nasv}
\end{align}
Equation~\eqref{eq:em-gbm-log} is the scheme implemented as \texttt{dX\_rhs} in
\texttt{SDEs/gbm.py}, and \eqref{eq:em-gbm} as \texttt{dS\_rhs}.
Equations~\eqref{eq:em-nasv} are implemented as \texttt{dXdY\_rhs} in
\texttt{SDEs/nasv.py}.

\subsubsection{Milstein}

The Milstein scheme~\cite{milstein} keeps the first stochastic correction that
EM discards. To derive it, expand the diffusion coefficient along the path,
\begin{equation}
    b(X_s) = b(X_n) + b'(X_n)(X_s - X_n)
             + \mathcal{O}\!\left( (X_s - X_n)^2 \right),
\end{equation}
and use the leading behaviour $X_s - X_n \approx b(X_n)(W_s - W_n)$, valid at
the order required. The It\^o integral in \eqref{eq:integral-form} then becomes
\begin{equation}
    \int_{t_n}^{t_{n+1}} b(X_s)\, dW_s
    \approx b(X_n)\, \Delta W_n
           + b(X_n) b'(X_n) \int_{t_n}^{t_{n+1}} (W_s - W_n)\, dW_s .
    \label{eq:milstein-integral}
\end{equation}
The remaining It\^o integral is elementary. Applying It\^o's lemma
\eqref{eq:ito} to $f(w) = \tfrac12 w^2$ gives
$d\!\left( \tfrac12 W_s^2 \right) = W_s\, dW_s + \tfrac12 ds$, hence
\begin{equation}
    \int_{t_n}^{t_{n+1}} (W_s - W_n)\, dW_s
    = \tfrac12 \left( \Delta W_n^2 - h \right).
    \label{eq:ito-key}
\end{equation}
Substituting \eqref{eq:ito-key} into \eqref{eq:milstein-integral} and neglecting
the higher-order terms gives the Milstein scheme
\begin{equation}
    X_{n+1} = X_n + a(X_n) h + b(X_n) \Delta W_n
              + \tfrac12 b(X_n) b'(X_n)
                \left( \Delta W_n^2 - h \right).
    \label{eq:milstein-general}
\end{equation}
The correction term is proportional to $b b'$, so it vanishes whenever the
diffusion coefficient is constant.

For GBM on the price, $a(S) = \mu S$, $b(S) = \sigma S$ and $b'(S) = \sigma$, so
$b b' = \sigma^2 S$ and
\begin{equation}
    S_{n+1} = S_n + \mu S_n h + \sigma S_n \Delta W_n
              + \tfrac12 \sigma^2 S_n \left( \Delta W_n^2 - h \right),
    \label{eq:milstein-gbm}
\end{equation}
which is \texttt{dS\_correction} in \texttt{SDEs/gbm.py}. For GBM on the
log-price, $b \equiv \sigma$ is constant, so the correction vanishes
identically: Milstein and Euler--Maruyama coincide on $X$. This is the precise
sense in which a strong order cannot be measured on the log-price scale, and it
is why \eqref{eq:milstein-gbm} rather than \eqref{eq:em-gbm-log} is used in the
order study of Section~\ref{sec:experiments}.

For the two-dimensional non-affine system, the same construction produces cross
terms and iterated stochastic integrals,
\begin{equation}
    z_{n+1} = z_n + a(z_n) h + B(z_n) \Delta \bm{W}_n
              + \sum_{j_1, j_2} L^{j_1} L^{j_2} b(z_n)\,
                I_{(j_1, j_2), n},
    \label{eq:milstein-multi}
\end{equation}
where $B$ collects the diffusion vector fields $b_1, b_2$ and
$I_{(j_1,j_2)}$ are iterated It\^o integrals. The terms with
$j_1 \ne j_2$ cannot be expressed through
$\Delta W^{(1)} \Delta W^{(2)}$ unless the vector fields commute. For the
non-affine model,
\begin{equation}
    b_1(X, Y) = \begin{pmatrix} g(Y) \\ 0 \end{pmatrix},
    \qquad
    b_2(X, Y) = \begin{pmatrix} 0 \\ \xi \sqrt{1 + Y^2} \end{pmatrix},
    \label{eq:vector-fields}
\end{equation}
and a direct computation gives
\begin{equation}
    [b_1, b_2] =
    \begin{pmatrix} -\xi \sqrt{1 + Y^2}\, g'(Y) \\ 0 \end{pmatrix} \ne 0,
    \label{eq:lie-bracket}
\end{equation}
because $g'(y) = (\sigma_{\max} - \sigma_{\min}) e^{-y} (1 + e^{-y})^{-2} > 0$
everywhere. The noise is therefore non-commutative, the cross terms do not
collapse, and a full Milstein method would require simulated L\'evy areas. The
scalar order-one result is not available for this system, so EM alone is used
for the non-affine model throughout.

\subsection{Consistency and Truncation Error}
\label{sec:truncation}

The Taylor expansion used above is incomplete, and the omitted terms set the
orders. For the scalar case, the It\^o--Taylor expansion to the first two
stochastic orders is
\begin{equation}
    X(t_{n+1}) = X_n + a h + b\, \Delta W_n
                 + \tfrac12 b b' \left( \Delta W_n^2 - h \right)
                 + R_n ,
    \label{eq:ito-taylor}
\end{equation}
where the remainder $R_n$ collects terms such as
$\int\!\int b\, a'\, dW\, ds$ and the triple It\^o integral.
Table~\ref{tab:ito-taylor} lists the leading terms with their size in the
$L^2(\Omega)$ norm and shows which scheme retains them.

\begin{table}[H]
\centering
\begin{tabular}{llcc}
\toprule
\textbf{It\^o--Taylor term} & \textbf{Size in $L^2(\Omega)$}
    & \textbf{EM} & \textbf{Milstein} \\
\midrule
$a\, h$                                & $\mathcal{O}(h)$       & \checkmark & \checkmark \\
$b\, \Delta W_n$                       & $\mathcal{O}(h^{1/2})$ & \checkmark & \checkmark \\
$\tfrac12 b b' \left( \Delta W_n^2 - h \right)$
                                       & $\mathcal{O}(h)$       & $\times$   & \checkmark \\
remainder $R_n$                        & $\mathcal{O}(h^{3/2})$ & $\times$   & $\times$ \\
\bottomrule
\end{tabular}
\caption{Leading It\^o--Taylor terms for a scalar SDE, their $L^2$ size, and
the schemes that retain them. Note that $\Delta W_n^2 - h$ has zero mean and
variance $2h^2$, so the third term is of size $\mathcal{O}(h)$, strictly larger
than the remainder.}
\label{tab:ito-taylor}
\end{table}

\paragraph{Euler--Maruyama}
The local defect of EM is the omitted third term plus the remainder,
\begin{equation}
    X(t_{n+1}) - X_{n+1}^{EM}
    = \underbrace{\tfrac12 b b' \left( \Delta W_n^2 - h \right)}
                 _{\mathcal{O}(h)}
    + \underbrace{R_n}_{\mathcal{O}(h^{3/2})},
    \label{eq:local-defect-em}
\end{equation}
so the local defect is of size $\mathcal{O}(h)$ in $L^2$.

\paragraph{Milstein}
Milstein retains that term, so its local defect is just the remainder,
\begin{equation}
    X(t_{n+1}) - X_{n+1}^{Mil} = R_n = \mathcal{O}\!\left( h^{3/2} \right).
    \label{eq:local-defect-mil}
\end{equation}

\paragraph{How the local defects accumulate}
This step differs from the deterministic theory and is worth stating carefully.
The leading local defect \eqref{eq:local-defect-em} has zero conditional
expectation and is uncorrelated across steps, because it depends only on
$\Delta W_n$. Defects of this kind do not add up linearly; they add up like the
increments of a random walk. Over $N = T/h$ steps the accumulated error is
therefore of size
\begin{equation}
    \sqrt{N} \times \mathcal{O}(h) = \mathcal{O}\!\left( h^{1/2} \right)
    \quad \text{for EM},
    \qquad
    \sqrt{N} \times \mathcal{O}(h^{3/2}) = \mathcal{O}(h)
    \quad \text{for Milstein}.
    \label{eq:accumulation}
\end{equation}
The factor $\sqrt{N}$ rather than $N$ is the reason why a strong order of $1/2$
appears even though every local step is accurate to $\mathcal{O}(h)$. This is
the main qualitative difference from the ODE case and it should not be smoothed
over.

Drift-type defects behave differently: they are not mean-zero and do accumulate
linearly, contributing $\mathcal{O}(h)$ after $N$ steps. This distinction is
what separates strong from weak convergence, and it is developed in
Section~\ref{sec:convergence} below.

\subsection{Stability Analysis}
\label{sec:ms-stability}

The problem pack replaces the deterministic stability and adaptivity analysis
of the ODE directions by convergence and uncertainty quantification for
Direction 4~\cite{kloeden,higham}. This subsection is therefore not required by
the brief. It is included for completeness and is correct, but its honest
conclusion is that no stability constraint binds in either model.

\subsubsection{Deterministic Reduction}

Setting $\sigma = 0$ in \eqref{eq:gbm-model} removes the noise. Both schemes
\eqref{eq:em-gbm} and \eqref{eq:milstein-gbm} then reduce to the explicit Euler
step $S_{n+1} = S_n (1 + \mu h)$, whose amplification factor is
\begin{equation}
    R(z) = 1 + z, \qquad z = \mu h,
    \label{eq:R-euler}
\end{equation}
with the classical absolute-stability region $|1 + z| < 1$. It should be noted
that this region is not informative here: $\mu = 0.05 > 0$, so the exact
solution grows and the notion of absolute stability, which describes decay,
does not apply. The benchmark GBM problem contains no stiff mode.

\subsubsection{Mean-Square Stability}

For SDEs the natural analogue is mean-square stability. Consider the linear
test equation $dX = \lambda X\, dt + \mu_d X\, dW$ with
$\operatorname{Re}\lambda < 0$, and define
\begin{equation}
    R_{MS}(z, w) = \frac{\mathbb{E}\!\left[ |X_{n+1}|^2 \right]}{|X_n|^2},
    \qquad z = \lambda h, \qquad w = \mu_d \sqrt{h} .
    \label{eq:ms-definition}
\end{equation}
The scheme is mean-square stable if $R_{MS}(z,w) < 1$; the exact solution has
$R_{MS}^{exact} = e^{2\lambda h + \mu_d^2 h}$.

\paragraph{Euler--Maruyama}
With $X_{n+1} = X_n (1 + \lambda h + \mu_d \Delta W_n)$ and
$\xi = \Delta W_n / \sqrt{h} \sim \mathcal{N}(0,1)$,
\[
    \mathbb{E}\!\left[ |1 + z + w \xi|^2 \right] = |1 + z|^2 + w^2,
\]
because $\mathbb{E}[\xi] = 0$ and $\mathbb{E}[\xi^2] = 1$. Hence
\begin{equation}
    R_{MS}^{EM}(z, w) = |1 + z|^2 + w^2 .
    \label{eq:ms-em}
\end{equation}
Mean-square stability requires $|1 + z|^2 + w^2 < 1$, that is the disc
$|1 + z| < \sqrt{1 - w^2}$, which is non-empty only if $w^2 < 1$, i.e.\
$\mu_d^2 h < 1$.

\paragraph{Milstein}
With the correction included,
$X_{n+1} = X_n \left[ 1 + z + w \xi + \tfrac12 w^2 (\xi^2 - 1) \right]$ and
using the Gaussian moments $\mathbb{E}[\xi^2] = 1$, $\mathbb{E}[\xi^3] = 0$ and
$\mathbb{E}[\xi^4] = 3$,
\begin{equation}
    R_{MS}^{Mil}(z, w) = |1 + z|^2 + w^2 + \tfrac12 w^4 .
    \label{eq:ms-mil}
\end{equation}
The extra term $\tfrac12 w^4$ is easily lost by copying \eqref{eq:ms-em}; its
presence matters when $w$ is not small. As a check, for $\lambda = 0$ the exact
factor is $e^{w^2} = 1 + w^2 + \tfrac12 w^4 + \mathcal{O}(w^6)$, which agrees
with \eqref{eq:ms-mil} to the order shown.

\paragraph{Application to GBM}
GBM in the price is linear with $\lambda = \mu$ and $\mu_d = \sigma$, so
$z = \mu h$ and $w = \sigma \sqrt{h}$. Since $\mu > 0$, the exact second moment
grows like $e^{(2\mu + \sigma^2) h}$, and $R_{MS} < 1$ cannot be attained by any
step size. Mean-square stability is therefore vacuous for this benchmark: there
is no decay for a scheme to preserve. What remains meaningful is the
\emph{accuracy} of the second moment. For EM,
\begin{equation}
    \mathbb{E}\!\left[ S_N^2 \right] = S_0^2
    \left[ (1 + \mu h)^2 + \sigma^2 h \right]^{N},
    \label{eq:second-moment-em}
\end{equation}
whereas the exact value is
$\mathbb{E}[S_T^2] = S_0^2 e^{(2\mu + \sigma^2) T}$. Comparing the two gives a
relative deviation of size
\begin{equation}
    \left( \mu^2 - \tfrac12 (2\mu + \sigma^2)^2 \right) T h
    + \mathcal{O}(h^2) = -0.0073\, h + \mathcal{O}(h^2)
    \label{eq:second-moment-bias}
\end{equation}
for the benchmark parameters. The deviation is of order $h$, which is exactly
the weak order reported in Section~\ref{sec:convergence}, and for the Milstein
scheme the same comparison improves by the term $\tfrac12\sigma^4h^2$.

\paragraph{The linearised non-affine system}
The Jacobian of the non-affine drift \eqref{eq:nasv-model-x}--%
\eqref{eq:nasv-model-y} is
\begin{equation}
    J(z) =
    \begin{pmatrix}
        0 & -g(Y) g'(Y) \\
        0 & -\kappa
    \end{pmatrix},
    \label{eq:jacobian-nasv}
\end{equation}
because the drift of $X$ depends on $Y$ alone and $X$ does not feed back. This
matrix is triangular, so its eigenvalues are $\{0, -\kappa\} = \{0, -2\}$ at
every state. The zero eigenvalue belongs to the neutral log-price direction,
and the only negative eigenvalue imposes the step restriction
$h < 2/\kappa = 1$, which is trivially satisfied by every step size used in
this report. As in the GBM case, stability is not the binding constraint;
accuracy is.

\subsection{Convergence}
\label{sec:convergence}

\paragraph{Why the deterministic route is not available}
For an ODE one proves convergence by combining consistency with a stability
estimate, and the Lax equivalence theorem closes the argument. No such theorem
exists for SDEs. The reason is that the local defects of a stochastic scheme
are random and their accumulation depends on their correlation structure, as
Section~\ref{sec:truncation} showed. The routes actually available are the
It\^o--Taylor expansion together with moment bounds, or the Gronwall argument
in Kloeden and Platen~\cite{kloeden}. The two orders are defined separately,
and they are not interchangeable.

\begin{definition}[Strong order]
\label{def:strong}
A scheme has strong order $\gamma$ if, for both the scheme and the exact
solution driven by the same Brownian path,
\begin{equation}
    \mathbb{E}\!\left[ \left| X_N - X(T) \right| \right]
    \le C h^{\gamma}
    \label{eq:strong-order}
\end{equation}
for all sufficiently small $h$. The error is a pathwise distance, so the two
solutions must be coupled through common random numbers.
\end{definition}

\begin{definition}[Weak order]
\label{def:weak}
A scheme has weak order $\beta$ if
\begin{equation}
    \left| \mathbb{E}\!\left[ f(X_N) \right]
           - \mathbb{E}\!\left[ f(X(T)) \right] \right|
    \le C h^{\beta}
    \label{eq:weak-order}
\end{equation}
for all $f$ in a suitable class of smooth test functions.
\end{definition}

\begin{theorem}[Convergence of EM and Milstein, scalar case]
\label{thm:convergence}
Under Assumption~\ref{ass:regularity}, Euler--Maruyama \eqref{eq:em-general}
has strong order $1/2$ and weak order $1$. If in addition $b \in C^1$ with
Lipschitz derivative, Milstein \eqref{eq:milstein-general} has strong order $1$
and weak order $1$.
\end{theorem}

Theorem~\ref{thm:convergence} is the reason for the two expectations used in
Section~\ref{sec:experiments}: the strong error \eqref{eq:strong-order} for EM
on the price is expected to fall like $h^{1/2}$, and for Milstein like $h$,
while the weak error \eqref{eq:weak-order} is expected to fall like $h$ for
both. Note that Milstein improves the strong order but not the weak one. The
added term has zero mean, so it does not change the bias; consequently the
expected value of $S_N$ is the same for EM and Milstein at every step size for
GBM. This is a sharp and easily tested prediction, and it is used as a
consistency check in Section~\ref{sec:experiments}.

\begin{remark}[The benchmark is not a test of the theorem]
\label{rem:exact-on-log}
On the log-price scale, EM reproduces the exact solution at grid points, so its
error is identically zero rather than $\mathcal{O}(h^{1/2})$.
Theorem~\ref{thm:convergence} is still consistent, but it is not sharp for the
degenerate reason discussed in Section~\ref{sec:schemes}. This is why the order
study is carried out on the price process \eqref{eq:em-gbm}, where the noise is
state-dependent and the theorem is informative.
\end{remark}

\begin{remark}[What can and cannot be measured for the non-affine model]
\label{rem:empirical}
For the non-affine system \eqref{eq:em-nasv} no exact solution exists, so the
strong error can only be computed against a finely resolved numerical solution.
Two consequences follow. First, the order obtained this way is an empirical
rate, not a theorem, and it must be reported as such. Second, because the noise
\eqref{eq:lie-bracket} is non-commutative, EM is expected to show a strong
order near $1/2$ and no scalar Milstein comparison is available. The reference
solution must itself be validated by refinement; agreement between two
discretisations is evidence, not an exact error calculation.
\end{remark}

\newpage
% ============================================================
\section{Implementation}
\label{sec:implementation}
% ============================================================

\subsection{Algorithms}
\label{sec:algorithms}

This section lists the algorithms that produce every number in
Section~\ref{sec:experiments}. Two schemes are used for GBM and one for the
non-affine model; the Brownian increments are produced once, outside the
schemes, by the engine of Section~\ref{sec:engine}, so that all step sizes and
both methods share the same underlying path.

\begin{algorithm}[H]
\caption{Euler--Maruyama (EM) for GBM, $dS = \mu S\,dt + \sigma S\,dW$}
\label{alg:em}
\begin{algorithmic}[1]
\Require initial price $S_0$, step size $h$, steps $N$,
         increments $\{\Delta W_n\}_{n=0}^{N-1}$, drift $\mu$, volatility $\sigma$
\Ensure approximate price $S_N$
\State $S \leftarrow S_0$
\For{$n = 0, 1, \ldots, N-1$}
    \State $S \leftarrow S + \mu\, S\, h + \sigma\, S\, \Delta W_n$
\EndFor
\State \Return $S$
\end{algorithmic}
\end{algorithm}

\begin{algorithm}[H]
\caption{Milstein for GBM, $dS = \mu S\,dt + \sigma S\,dW$}
\label{alg:milstein}
\begin{algorithmic}[1]
\Require initial price $S_0$, step size $h$, steps $N$,
         increments $\{\Delta W_n\}_{n=0}^{N-1}$, drift $\mu$, volatility $\sigma$
\Ensure approximate price $S_N$
\State $S \leftarrow S_0$
\For{$n = 0, 1, \ldots, N-1$}
    \State $S \leftarrow S + \mu\, S\, h + \sigma\, S\, \Delta W_n
                    + \tfrac12\, \sigma^2 S \left( \Delta W_n^2 - h \right)$
\EndFor
\State \Return $S$
\end{algorithmic}
\end{algorithm}

\begin{algorithm}[H]
\caption{Euler--Maruyama for the non-affine model, state $z_n = (X_n, Y_n)$}
\label{alg:em-nasv}
\begin{algorithmic}[1]
\Require input $(X_0, Y_0)$, step size $h$, steps $N$,
         correlated increments $\{\Delta W_n^{(1)}, \Delta W_n^{(2)}\}$,
         parameters $\mu, \kappa, \theta, \xi, \sigma_{\min}, \sigma_{\max}$
\Ensure $X_N$; the price is $S_N = e^{X_N}$
\State $X \leftarrow X_0$, \; $Y \leftarrow Y_0$
\For{$n = 0, 1, \ldots, N-1$}
    \State $\sigma_n \leftarrow g(Y)$
           \Comment{logistic volatility, \eqref{eq:nasv-g}}
    \State $X \leftarrow X + \left( \mu - \tfrac12 \sigma_n^2 \right) h
              + \sigma_n\, \Delta W_n^{(1)}$
    \State $Y \leftarrow Y + \kappa (\theta - Y)\, h
              + \xi \sqrt{1 + Y^2}\; \Delta W_n^{(2)}$
\EndFor
\State \Return $X$
\end{algorithmic}
\end{algorithm}

No Milstein counterpart to Algorithm~\ref{alg:em-nasv} is implemented. As shown
in \eqref{eq:lie-bracket}, the two diffusion vector fields of the non-affine
system do not commute, so a multidimensional Milstein method would require
cross terms and simulated L\'evy areas. The solver therefore raises
\texttt{NotImplementedError} if it is asked for a Milstein step on this model.
This is deliberate: it makes it impossible to produce an order-one claim for
the non-affine system by accident.

\subsection{Brownian Engine and Shared Paths}
\label{sec:engine}

Strong error is a pathwise distance (Definition~\ref{def:strong}), so the two
solutions being compared must be driven by the same Brownian path. If instead a
fresh path were drawn for each step size, the differences between levels would
be dominated by sampling noise rather than by the discretisation error, and no
order could be estimated. The engine therefore works as follows.

Fine increments are generated once at the finest resolution $n_{\max}$, and the
increments at any coarser resolution $n$ with $n \mid n_{\max}$ are obtained by
summing blocks of fine increments:
\begin{equation}
    \left( \Delta W^{(c)} \right)_j
    = \sum_{i = j b + 1}^{(j+1) b} \Delta W_i,
    \qquad b = n_{\max} / n, \qquad j = 0, \ldots, n-1 .
    \label{eq:coarsening}
\end{equation}
Because the coarse increments are block sums of the fine increments, the coarse
path is exactly the fine path sampled on the coarse grid:
\begin{equation}
    \sum_{j} \left( \Delta W^{(c)} \right)_j = \sum_{i} \Delta W_i = W(T)
    \qquad \text{for every } n .
    \label{eq:same-endpoint}
\end{equation}
Equation~\eqref{eq:same-endpoint} is what makes the strong error comparable
across levels: the reference value $S(T)$ against which each scheme is measured
is the same random variable at every resolution.

Two further properties of \eqref{eq:coarsening} matter. First, the correlation
is preserved exactly: the coarse pair is still jointly Gaussian with
$\operatorname{Cov}\!\left( \Delta W^{(1)}, \Delta W^{(2)} \right)
= \rho\, \Delta t$, so the parameter $\rho$ is constant across the whole study
even though the increments are never regenerated. Second, the nesting is
transitive, so the levels can be produced by successive halving of the finest
log. The implementation keeps only the current level in memory and discards the
previous one, which bounds the memory by roughly twice the finest level rather
than by the sum over all levels.

\subsection{Code Structure}
\label{sec:code-structure}

\begin{verbatim}
code/
|-- run_all.py                   # runs every experiment, writes data/ and figures/
|-- SDEs/
|   |-- bm_engine.py             # standard_bm, correlated_bm_pair, extract_dw
|   |-- gbm.py                   # dS_rhs, dS_correction, dX_rhs, S_exact
|   `-- nasv.py                  # g, dXdY_rhs, benchmark parameters
|-- solvers/
|   |-- em.py                    # em_step, em_solve           (Euler-Maruyama)
|   `-- milstein.py              # milstein_step, milstein_solve (scalar GBM only)
|-- tools/
|   |-- mc.py                    # Brownian simulator + increment-log I/O
|   `-- stats.py                 # mean, standard error, CI, log-log slope fit
|-- experiments/
|   |-- exp0_simulate_bm.py      # shared-path smoke test
|   |-- exp1_strong_gbm.py       # strong convergence, EM vs Milstein
|   |-- exp2_weak_gbm.py         # weak convergence and the closed-form bias
|   |-- exp3_distribution.py     # distribution of S(T) vs the lognormal law
|   |-- exp4_positivity.py       # negativity rate of EM and Milstein
|   |-- exp5_strong_nasv.py      # coupled convergence against fine references
|   |-- exp6_payoff_nasv.py      # weak convergence of E[(S_T - 100)^+]
|   `-- exp7_diagnostics.py      # increment moments, finiteness, positivity
|-- data/                        # increment logs and result tables (CSV/JSON)
`-- figures/                     # figures used in this report
\end{verbatim}

The two model modules contain only definitions; the two solver modules contain
only stepping; the experiments contain only assembly and reporting. No
experiment re-implements a scheme or draws its own noise. The solver state for
the non-affine model is the pair $(X, Y)$ rather than $(S, Y)$, so the price is
always recovered as $S = e^{X}$ and positivity is structural rather than
enforced by a projection.

\subsection{Notation in the Code}
\label{sec:notation-code}

Table~\ref{tab:notation-code} maps the symbols used in this report onto the
names used in the code, so that every formula can be located in the repository
without ambiguity.

\begin{table}[H]
\centering
\small
\begin{tabular}{lll}
\toprule
\textbf{Report symbol} & \textbf{Code name} & \textbf{Location} \\
\midrule
$S_0$                         & \texttt{GBM\_S0}, \texttt{NA\_SV\_S0}
                              & \texttt{SDEs/gbm.py}, \texttt{SDEs/nasv.py} \\
$\mu, \sigma$                 & \texttt{GBM\_MU}, \texttt{GBM\_SIGMA}
                              & \texttt{SDEs/gbm.py} \\
$\kappa, \theta, \xi, \rho$   & \texttt{NA\_SV\_KAPPA}, \texttt{NA\_SV\_THETA},
                                \texttt{NA\_SV\_XI}, \texttt{NA\_SV\_RHO}
                              & \texttt{SDEs/nasv.py} \\
$\sigma_{\min}, \sigma_{\max}$ & \texttt{NA\_SV\_SIGMA\_MIN},
                                 \texttt{NA\_SV\_SIGMA\_MAX}
                              & \texttt{SDEs/nasv.py} \\
$g(\cdot)$                    & \texttt{g} & \texttt{SDEs/nasv.py} \\
$h$                           & \texttt{dt}
                              & \texttt{em\_solve}, \texttt{milstein\_solve} \\
$N$                           & \texttt{n\_steps}
                              & \texttt{em\_solve}, \texttt{milstein\_solve} \\
$n_{\max}$                    & \texttt{n\_max}
                              & \texttt{bm\_engine}, \texttt{mc} \\
$\Delta W_n$                  & \texttt{dW} & \texttt{bm\_engine} \\
$\Delta W_n^{(1)}, \Delta W_n^{(2)}$ & \texttt{dW1}, \texttt{dW2}
                              & \texttt{bm\_engine} \\
$z_n = (X_n, Y_n)$            & \texttt{X}, \texttt{Y} & \texttt{solvers/em.py} \\
$S_N, X_N$                    & \texttt{out[:, -1]}
                              & \texttt{em\_solve}, \texttt{milstein\_solve} \\
\bottomrule
\end{tabular}
\caption{Report notation mapped to code identifiers.}
\label{tab:notation-code}
\end{table}

\subsection{Key Design Decisions}
\label{sec:key-decisions}

\begin{itemize}
\item \textbf{Shared Brownian path.} As explained in Section~\ref{sec:engine},
      the engine draws fine increments once at $n_{\max}$ and coarsens by
      summing blocks, so every step size and both methods see the same
      realisation. The increment logs are written to \texttt{data/} and
      reloaded by each experiment, so the coupling holds \emph{between}
      experiments as well as within one. The cost is storage of order
      $2 n_{\max} M$ doubles for the correlated pair, which is the reason
      Algorithm~\ref{alg:em-nasv} is run in batches over paths rather than on
      the full ensemble at once.
\item \textbf{Terminal values rather than trajectories.} The convergence and
      payoff estimators only need $S_N$, not the whole path, so the solvers are
      called in a mode that stores the current state and, where a trajectory is
      required for a figure, only for the small cases. This keeps the memory
      footprint at $\mathcal{O}(M)$ per level instead of
      $\mathcal{O}(M n_{\max})$.
\item \textbf{Vectorisation over paths, not over time.} The stepping loops run
      over the $N$ time steps and operate on the whole path ensemble at once.
      For GBM this is not strictly necessary, because \eqref{eq:em-gbm} and
      \eqref{eq:milstein-gbm} can be written as cumulative products of per-step
      factors; this shortcut is used where the terminal value is all that is
      required. For the non-affine model it is necessary, because the
      coefficients of step $n+1$ depend on the state reached at step $n$ and no
      cumulative representation exists.
\item \textbf{Log-price transform.} Discretising $X = \log S$ makes the GBM
      noise additive, so the exact sampler \eqref{eq:gbm-sampler} needs no
      discretisation at all, and $S = e^{X} > 0$ holds automatically. For the
      non-affine model the same substitution keeps the price positive without a
      projection step, which would otherwise bias the weak error.
\item \textbf{Milstein only for scalar GBM.} The correction
      \eqref{eq:milstein-gbm} is implemented for the price process, where
      $b' \ne 0$. It is not implemented for the log-price (where it vanishes)
      or for the non-affine system (where the noise is non-commutative,
      \eqref{eq:lie-bracket}). The solver raises an error rather than returning
      a silently wrong order.
\item \textbf{A separate exact reference.} The exact solution
      \eqref{eq:gbm-exact} is a distinct function, not a scheme, and it is
      never listed among the methods. This prevents the exact sampler from
      being reported as a numerical method with zero error.
\item \textbf{Every estimate carries a standard error.} All Monte Carlo
      quantities are reported as
      $\hat{\theta} \pm 1.96\, \widehat{\mathrm{se}}$, and a log--log slope is
      quoted only when the estimated difference between adjacent levels exceeds
      twice its paired standard error. Because the coarse and fine paths are
      coupled, that difference is computed on paired samples rather than on two
      independent batches.
\item \textbf{Numerical evaluation of the logistic function.} The volatility
      \eqref{eq:nasv-g} is evaluated as $1/(1 + e^{-y})$, which overflows for
      large negative $y$. The implementation uses a branch (or \texttt{expit})
      that returns $\sigma_{\min}$ in that regime. Under the benchmark
      parameters $Y$ remains of order one, so this is a robustness safeguard
      rather than an active concern.
\end{itemize}

\subsection{Implementation Verification}
\label{sec:impl-verify}

Before any convergence result is quoted, the following checks are run. They are
cheap, they are independent of each other, and each one catches a different
class of error.

\begin{enumerate}[leftmargin=2em]
    \item \textbf{Increment moments.} For the correlated pair, the sample
          variance of $\Delta W_n^{(1)}$ and $\Delta W_n^{(2)}$ should agree
          with $h$, and their sample correlation with $\rho$. This detects an
          error in the Cholesky step \eqref{eq:increments}.
    \item \textbf{Coupling identity.} For every $n \mid n_{\max}$, the
          cumulative sum of the coarse increments must equal the fine
          cumulative sum at the coarse grid points, to round-off. This detects
          an error in the block size or in the reshaping of the increment log.
    \item \textbf{Exact moments for GBM.} The sample mean and variance of the
          exact sampler are compared with \eqref{eq:gbm-moments}. For
          $M = 10^5$ paths the standard error of the mean is $0.067$, so the
          sample mean must fall within about $0.13$ of $105.1271$.
    \item \textbf{Exact mean for the non-affine model.} As shown in
          Remark~\ref{rem:nasv-mean}, $\mathbb{E}[S_T] = 105.1271$ holds
          independently of every volatility parameter. This is the only exact
          identity available for that model and it is checked at every step
          size, because a sign or factor error in the $-\tfrac12 g(Y)^2$ term
          shows up immediately as a shift in the sample mean.
    \item \textbf{Equality of the two weak curves.} By
          Theorem~\ref{thm:convergence} and the zero mean of the Milstein
          correction, EM and Milstein must give the same $\mathbb{E}[S_N]$ for
          GBM at every step size. Any visible gap between the two
          weak-convergence curves indicates an implementation error, not a
          discretisation effect.
    \item \textbf{Finiteness and positivity.} All simulated paths must be finite
          for both models, and $S = e^{X} > 0$ must hold throughout for the
          non-affine model.
    \item \textbf{Reference resolution.} The numerical reference for the
          non-affine model is computed at two successive refinements, and the
          discrepancy between them must be smaller than the quantities being
          measured. Without this check the empirical rate of
          Remark~\ref{rem:empirical} cannot be interpreted.
\end{enumerate}

\subsection{Reproducibility}
\label{sec:reproducibility}

Each experiment fixes its own generator seed, and the increment logs in
\texttt{data/} are stored together with the seeds and parameters that produced
them, so that every figure and table can be regenerated from
\texttt{run\_all.py} alone. Table~\ref{tab:experiment-design} records the design
of each experiment: the model, the methods compared, the step sizes, the number
of paths and the cost metrics that are reported.

\begin{table}[H]
\centering
\small
\begin{tabular}{lllll}
\toprule
\textbf{Script} & \textbf{Model} & \textbf{Methods} & \textbf{Step sizes}
    & \textbf{Paths} \\
\midrule
\texttt{exp1} & GBM        & EM, Milstein, exact
              & $h = T/2^k$, $k = 4, \ldots, 12$ & $10^5$ \\
\texttt{exp2} & GBM        & EM, Milstein
              & same                            & $10^5$ \\
\texttt{exp3} & GBM        & exact, EM
              & $h = 1/252$, $1/1024$           & $10^5$ \\
\texttt{exp4} & GBM        & EM, Milstein
              & $h = 1/2^k$, $k = 2, \ldots, 8$ & $10^6$ \\
\texttt{exp5} & non-affine & EM
              & $h = T/2^k$, $k = 4, \ldots, 10$ & $10^4$ \\
\texttt{exp6} & non-affine & EM
              & $h = T/2^k$, $k = 3, \ldots, 9$  & $10^4$ \\
\texttt{exp7} & both       & ---
              & ---                              & $10^4$ \\
\bottomrule
\end{tabular}
\caption{Design of the numerical experiments. The non-affine experiments use
fine references at $h_{\mathrm{ref}} = T/2^{13}$ and $T/2^{14}$; each experiment
records the number of paths, the number of time steps, the wall time and the
number of random numbers consumed.}
\label{tab:experiment-design}
\end{table}

\newpage
% ============================================================
\section{Numerical Experiments}
\label{sec:experiments}
% ============================================================

\subsection{Verification}
\label{sec:verification}

\subsubsection{GBM with Exact Solution}

[Present the exact-solution test. Report the sample mean and variance of the
exact sampler against \eqref{eq:gbm-moments}, with standard errors, and the
moment checks of Section~\ref{sec:impl-verify}. This establishes that the
engine, the reference and the estimators agree before any order is quoted.]

\subsubsection{Non-Affine Model: Numerical Reference and Nested-Grid Validation}

[The reference here is a finely resolved numerical solution, not a manufactured
solution, and it must be described as such. Report the reference at
$h_{\mathrm{ref}} = T/2^{13}$ and $T/2^{14}$, the discrepancy between the two,
and the checks of Section~\ref{sec:impl-verify}. State explicitly that
agreement between two discretisations is evidence, not an exact error
calculation.]

\begin{table}[H]
\centering
\begin{tabular}{lcccc}
\toprule
\textbf{Step size $h$} & \textbf{Error} & \textbf{Std.\ error}
    & \textbf{Ratio} & \textbf{Order} \\
\midrule
$T/2^{4}$   & [value] & [value] & ---     & ---     \\
$T/2^{5}$   & [value] & [value] & [ratio] & [order] \\
$T/2^{6}$   & [value] & [value] & [ratio] & [order] \\
$T/2^{7}$   & [value] & [value] & [ratio] & [order] \\
$T/2^{8}$   & [value] & [value] & [ratio] & [order] \\
\bottomrule
\end{tabular}
\caption{Convergence study for [test case], with $M$ paths per level. A slope is
quoted only for the levels where the difference between successive errors
exceeds twice its paired standard error.}
\label{tab:convergence}
\end{table}

\subsection{Strong Convergence}
\label{sec:strong}

[Plot $\mathbb{E}|S_N - S(T)|$ against $h$ on log--log axes for EM and Milstein
on the price process, with error bars. Expected slopes: $1/2$ for EM and $1$ for
Milstein. Report the fitted slopes with confidence intervals and state the range
of $h$ over which each fit was made.]

\subsection{Weak Convergence and Monte Carlo Uncertainty}
\label{sec:weak}

[Report three separate pieces of evidence: the closed-form bias
$\mathbb{E}[S_N] = S_0 (1 + \mu h)^{T/h}$ compared with $S_0 e^{\mu T}$; a Monte
Carlo estimate using common random numbers with confidence intervals; and the
predicted equality of the EM and Milstein weak curves from
Theorem~\ref{thm:convergence}. For the payoff $\mathbb{E}[(S_T - 100)^+]$, use
the numerical reference with the same random numbers and report a slope only
when the estimated bias exceeds its sampling uncertainty.]

\subsection{Distribution of $S(T)$ and Positivity}
\label{sec:distribution}

[Compare the simulated distribution of $S(T)$ with the lognormal law of
Section~\ref{sec:gbm-model} by histogram overlay and quantile plot. For the
positive-preservation study, report the observed negativity rate of the EM and
Milstein updates in the price against $h$, together with the analytic bound:
the EM bracket $1 + \mu h + \sigma \Delta W_n$ is negative exactly when
$\Delta W_n < -(1 + \mu h)/\sigma$, an event of probability
$\Phi\!\left( -(1 + \mu h)/(\sigma\sqrt{h}) \right) \sim e^{-1/(2\sigma^2 h)}$,
while the Milstein bracket has minimum $\tfrac12 + (\mu - \tfrac12\sigma^2)h$, so
with the parameters of Table~\ref{tab:gbm-params} it cannot become negative for
any $h$. State which parameter values are needed to make the effect visible and
that these are a stress test rather than a realistic setting.]

\subsection{Cost versus Accuracy}
\label{sec:cost}

[Report paths, time steps, wall time and random numbers consumed for each method
and step size. Note that EM and Milstein consume the same random numbers, so the
comparison is like for like.]

\begin{figure}[H]
\centering
\fbox{\parbox{0.8\textwidth}{\centering \vspace{3em}
[Insert your figure here] \vspace{3em}}}
\caption{[Caption that interprets rather than describes. E.g.\ ``The log--log
plot gives slopes of $0.49$ for Euler--Maruyama and $0.98$ for Milstein, with
the Milstein curve flattening above $h \approx 2^{-10}$ where the reference
error becomes visible.'']}
\label{fig:convergence}
\end{figure}

\subsection{Limitations and Robustness}
\label{sec:limitations}

[Discuss the regimes where the conclusions weaken: very coarse steps, where the
non-affine reference is insufficiently resolved; parameter values that make the
correlated increments poorly conditioned; and the fact that no Milstein
comparison exists for the non-affine model. This is not a weakness; it is
scientific honesty.]

\newpage
% ============================================================
\section{Conclusions}
\label{sec:conclusions}
% ============================================================

[Summarise in 3--4 bullet points:]
\begin{itemize}[leftmargin=2em]
\item [What you achieved.]
\item [The most surprising or instructive finding. Candidates: that the two
      weak curves coincide exactly, that the Milstein second-moment factor
      carries an extra $\tfrac12 w^4$ term, or that no stability constraint
      binds in either model.]
\item [The main limitation you identified: that the non-affine reference is
      numerical, so its rate is empirical.]
\item [A concrete direction for future work: for example, replacing EM on the
      non-affine system by a method with simulated L\'evy areas, or calibrating
      the model to quoted option prices.]
\end{itemize}

\newpage
% ============================================================
\section*{References}
% ============================================================

\begin{thebibliography}{9}

\bibitem{blackscholes} F. Black and M. Scholes, ``The pricing of options and
    corporate liabilities,'' \emph{Journal of Political Economy} 81(3):637--654,
    1973.

\bibitem{milstein} G. N. Milstein, ``Approximate integration of stochastic
    differential equations,'' \emph{Theory of Probability and its Applications}
    19(3):557--562, 1974.

\bibitem{kloeden} P. E. Kloeden and E. Platen, \emph{Numerical Solution of
    Stochastic Differential Equations}. Springer, 1992.

\bibitem{higham} D. J. Higham, ``An algorithmic introduction to numerical
    simulation of stochastic differential equations,'' \emph{SIAM Review}
    43(3):525--546, 2001.

\bibitem{dfs} D. Duffie, D. Filipovi\'c and W. Schachermayer, ``Affine processes
    and applications in finance,'' \emph{Annals of Applied Probability}
    13(3):984--1053, 2003.

\bibitem{heston} S. L. Heston, ``A closed-form solution for options with
    stochastic volatility with applications to bond and currency options,''
    \emph{Review of Financial Studies} 6(2):327--343, 1993.

\bibitem{numpy} C. R. Harris, K. J. Millman, S. J. van der Walt et al.,
    ``Array programming with NumPy,'' \emph{Nature} 585:357--362, 2020.

\end{thebibliography}

\newpage
% ============================================================
\appendix
\section{AI Transparency Log}
% ============================================================

[Append the completed team AI Transparency Log here. See
\texttt{06\_Templates/ICS} for the full form: tools used, significant claim
audits, what you changed and why, and the intellectual ownership declarations.]

\section{Individual Contribution Statements}
% ============================================================

Each team member must complete the following (repeat for all 5 members):

\begin{formbox}
\textbf{Name:} \underline{\hspace{5cm}}

\vspace{0.5em}
\textbf{1. My specific contributions to this project:}

[3--5 sentences: which equations you derived, which code modules you wrote,
which figures you produced, which report sections you drafted.]

\vspace{0.5em}
\textbf{2. The hardest conceptual obstacle I encountered:}

[One specific mathematical or algorithmic difficulty and how you overcame it.]

\vspace{0.5em}
\textbf{3. One piece of AI-generated output my team used, and what I changed:}

[Cross-reference an item from the AI Transparency Log.]

\vspace{0.5em}
\textbf{4. One thing I still do not fully understand:}

[Not penalised; helps the instructor plan the next project.]

\vspace{0.5em}
\textbf{Signature:} \underline{\hspace{4cm}} \quad
\textbf{Date:} \underline{\hspace{2cm}}
\end{formbox}

\section{Supplementary Figures and Data}
% ============================================================

[Place any additional figures, tables, or code listings that would interrupt
the main narrative but may be useful to a curious reader.]

\end{document}